\documentclass{article}
\usepackage{graphicx} % Required for inserting images
\usepackage{amsmath}
\usepackage{amssymb} 
\usepackage[authoryear]{natbib}
\bibliographystyle{apalike}
\usepackage[font=small,labelfont=bf]{caption}
\usepackage{authblk}
\usepackage[hidelinks]{hyperref}



\title{Single-molecule multimodal timing of in vivo mRNA synthesis}
\author{JC notes}

\date{24th March, 2026}

\begin{document}	
	
	\maketitle
	
	\section{Abstract}
	\begin{itemize}
		\item Poly(A) tail measurement to distinguish transcribing vs. poly(A) transcripts.
		\item Splicing is much more delayed than previously tought, for at least 10 kb. 
		\item Splicing mostly apparent only in poly(A) transcripts $\to$ they think it's delayed mostly after 3' end formation. 
		\item Lots of $\text{m}^{6}$A on unspliced RNA $\to$ deposition of $\text{m}^{6}$A happens before splicing.
		\item $\text{m}^{6}$A supressed around exon boundaries $\to$ $\text{m}^{6}$A topology should be established before EJC deposition.
		\item 3' end cleavage occurs multiple kbs behind transcription and is coordinated with transcription termination. The 3' end formation is recursive in many genes.
		\item Overall, they think that cleavage, termination and $\text{m}^{6}$A deposition occur earlier than most of the splicing.
	\end{itemize}
	
	\section{Intro}
	\begin{itemize}
		\item Spliceosome assembly is known to be rapid. 
		\item Splicing kinetics: estimates vary across methods.
		\item Previously, it was tought that EJC prevents $\text{m}^{6}$A deposition at exon boundaries. Here, they argue that $\text{m}^{6}$A is placed before splicing (and thus, before EJC placement).
		\begin{itemize}
			\item $\text{m}^{6}$A methylation machinery is known to interact with Pol II $\to$ permitting co-transcriptional methylation.
		\end{itemize}
		
	\end{itemize}
	
	\section*{Poly(A) notes}
	\begin{itemize}
		\item Poly(A) site: cannonical AAUAAA motif.
		\item When Pol II reaches it, it keeps transcribing.
		\item Protein complex cut RNA there; poly(A) polymerase (PAP) adds ~ 200 As.
		\item Transcription termination:
		\begin{itemize}
				\item Torpedo model: Exonuclease XRN2 start chewing the 5' end of RNA still transcribed by Pol II $\to$ when it catches Pol II, it knocks it of DNA.
				\item Allosteric model: Pol II changes conformation at poly(A) site, becomes less stable and eventually falls of DNA.
		\end{itemize}
	\end{itemize}
	
	\section{Results}
	\subsection{Development of a multimodal direct pre-mRNA sequencing technique}
	\begin{itemize}
		\item Need to distinguish whether pre-mRNA has a poly(A) tail:
		\begin{itemize}
			\item pre-mRNA\textsuperscript{e}: undergoing transcription.
			\item pre-mRNA\textsuperscript{a}: polyadenylated pre-mRNA, undergoing post-transcriptional maturation.
		\end{itemize}
		\item They pulled Pol II with associated pre-mRNA via immunoprecipitation. The pre-mRNA\textsuperscript{a} was still somehow pulled (even thought it was already cleaved - otherwise it would not have the poly(A) tail); probably was somehow attached via the cleavage complex, or some chromatin components.
		\item They got about 90\% of pre-mRNA\textsuperscript{e} and 10\% pre-mRNA\textsuperscript{a}.  
		\item 95\% of pre-mRNA\textsuperscript{e} have unspliced introns, while this is only 58\% for  pre-mRNA\textsuperscript{a}.
		\item Nanopore sequencing requires poly(A) tail $\to$ added artificially during library preparation. This \textit{in vitro} poly(A) tail is shorter (20-80nt), so they get bimodal distribution of tail lengths, and classify by cut-off.
		\begin{itemize}
			\item Also, the reads classified as pre-mRNA\textsuperscript{a} have the start of poly(A) highly enriched by known poly(A) sites.
		\end{itemize}
		
	\end{itemize}

	
	
	
	
	
	
	
	
	
	
	
	
	
	
\end{document}
